The current problem setting is evidence-grounded extractive question answering
with a constrained binary action:
\begin{equation}
f(q, e) \rightarrow \{\text{answer span}, \abstain\},
\end{equation}
where $q$ is a question and $e$ is a single evidence passage.

The Stage 1 scope is intentionally strict:
\begin{itemize}[leftmargin=1.25em]
  \item Dataset: SQuAD~2.0 only.
  \item Input: one question and one evidence passage.
  \item Output: extractive answer span or \abstain.
  \item Current implementation goal: reduce unsupported answers without sharply
  degrading performance on answerable questions.
\end{itemize}

This scope excludes retrieval, multi-hop reasoning, custom architectures,
long-form generation, and broader claim-level verification.
Those extensions are part of later roadmap stages, but they are out of scope
for the current paper draft's core experimental claim.

We evaluate the task with metrics that separate utility from caution.
For answerable examples, we track exact match and token-level F1.
For unanswerable examples, we track unsupported answer rate and abstention
quality.
This makes it possible to study the trade-off between being useful and being
appropriately cautious instead of collapsing everything into a single accuracy
number.
